\documentclass{article}

% if you need to pass options to natbib, use, e.g.:
%     \PassOptionsToPackage{numbers, compress}{natbib}
% before loading agents4science_2025

% ready for submission
\usepackage{agents4science_2025}

% to compile a preprint version, e.g., for submission to arXiv
% ...preamble continues...
% [preprint] option:
%     \usepackage[preprint]{agents4science_2025}

% to compile a camera-ready version, add the [final] option, e.g.:
%     \usepackage[final]{agents4science_2025}

% to avoid loading the natbib package, add option nonatbib:
%    \usepackage[nonatbib]{agents4science_2025}

% For workshops, the authors should use the workshop options and add the name of the workshop. 
% The "\workshoptitle" command is used to set the workshop title.
%
% \usepackage[sglblindworkshop]{agents4science_2025}
% \workshoptitle{WORKSHOP TITLE}


\usepackage[utf8]{inputenc} % allow utf-8 input
\usepackage[T1]{fontenc}    % use 8-bit T1 fonts
\usepackage{hyperref}       % hyperlinks
\usepackage{url}            % simple URL typesetting
\usepackage{booktabs}       % professional-quality tables
\usepackage{amsfonts}       % blackboard math symbols
\usepackage{nicefrac}       % compact symbols for 1/2, etc.
\usepackage{microtype}      % microtypography
\usepackage{xcolor}         % colors
\usepackage{graphicx}      % figures
\usepackage[strings]{underscore} % allow underscores in verbatim-like text (CSV import)
\usepackage{csvsimple}    % CSV/TSV tables
% \usepackage{verbatim}     % optional raw inclusion

% Note: For the workshop paper template, both \title{} and \workshoptitle{} are required.
% \title{} is the paper title and \workshoptitle{} is the workshop title for the footnote. 
  	\title{Mutual Wanting in Human--AI Interaction: A Minimal Method Study of User and System Desiderata Across GPT Model Generations}


% The \author macro works with any number of authors. There are two commands
% used to separate the names and addresses of multiple authors: \And and \AND.
%
% Using \And between authors leaves it to LaTeX to determine where to break the
% lines. Using \AND forces a line break at that point. So, if LaTeX puts 3 of 4
% authors names on the first line, and the last on the second line, try using
% \AND instead of \And before the third author name.


\author{Anonymous Authors\\Submission Under Review}


\begin{document}


\maketitle


\begin{abstract}
The rapid, often controversial, evolution of large language models (LLMs) necessitates a deeper understanding of the user-AI relationship. We introduce the concept of "mutual wanting" to analyze the bidirectional expectations between users and LLM agents, particularly during major version transitions. Using a minimal-resource approach that combines large-scale analysis of Reddit discourse with controlled API probing of different GPT versions, we investigate what users want from AI personas and what behaviors AI systems implicitly "want" from users. We propose a conceptual framework, the Mutual Wanting Alignment Framework (M-WAF), to map these desires along dimensions of uncertainty disclosure, relational pacing, and intimacy boundaries. This paper outlines our methodology, specifies our analytical pipeline for identifying tensions (e.g., warmth vs. cost-efficiency, stability vs. iteration), and details our planned metrics and ethical guardrails for designing more stable and trustworthy relational AI.\end{abstract}


% (Template instructional section removed for clarity; manuscript now begins directly with Introduction.)
\section{Introduction}
The rapid evolution of large language models (LLMs) has led to increasingly complex and often fraught interactions between users and AI systems. Large-scale model upgrades, such as the recent transition from GPT-4 to GPT-5, frequently trigger significant affective and epistemic friction. Public forums like Reddit are replete with user discourse expressing perceived losses in model "persona," such as diminished warmth, reduced creativity, or shifts in helpfulness. We conceptualize these reactions not as mere preference volatility, but as signals of a structural misalignment between user desires and system capabilities.

This paper introduces the concept of \emph{mutual wanting} to describe this bidirectional relationship. On one hand, users have explicit and implicit wants regarding the AI's relational, epistemic, and agentic affordances. On the other, the AI system, through its design and optimization, implicitly "wants" certain user behaviors, such as providing clear, structured inputs. We argue that mapping and reconciling these wants is a critical challenge for human-AI interaction.

To this end, we propose a minimal-resource research agenda that leverages publicly available data to explore this dynamic. By combining large-scale analysis of Reddit discourse with controlled, comparative API probing of different GPT versions, we can illuminate generalizable principles for designing more stable and trustworthy relational AI. This paper outlines our conceptual framework, details our analytical pipeline, and specifies the metrics and ethical guardrails for this investigation.

\paragraph{Contributions}
This paper makes several contributions. First, it introduces the Mutual Wanting Alignment Framework (M-WAF), a conceptual tool for deconstructing the bidirectional relationship between users and AI systems. The framework identifies key dimensions of user wants (e.g., epistemic reliability, relational continuity) and system wants (e.g., feedback efficiency, emotional scope boundaries). Second, we outline a multi-layer pipeline for analyzing public discourse, combining topic modeling, stance clustering, and persona drift tracking. Third, we propose a comparative behavioral probing protocol to quantify inter-version differences in how models handle uncertainty, silence, and intimacy. Finally, we specify a suite of metrics for measuring trust calibration and relational tension, and provide a draft of ethical guidelines for designing and deploying evolving AI personas.

\section{Related Work}
Our study builds on several streams of research in human-AI interaction. A central theme is the management of user perceptions and trust as AI systems evolve. Guidelines for human-AI interaction emphasize the importance of setting clear expectations and communicating capability changes to maintain user trust \citep{amershi2019guidelines}. Research in trust calibration shows that system transparency, for instance through explanations or uncertainty disclosures, is crucial for aligning user confidence with system capabilities \citep{yin2019understanding, lai2019human}. The way uncertainty is presented can significantly influence user trust and reliance on AI predictions \citep{kay2016uncertainty, reyes2025uncertainty}.

Another critical area is the socio-emotional dimension of human-AI relationships. The tendency for users to anthropomorphize AI agents is a well-documented phenomenon, driven by motivations for social connection and understanding \citep{epley2007seeing}. This can lead to the formation of parasocial relationships, where users develop one-sided bonds with AI personas \citep{horton1956mass}. While such relationships can enhance engagement, they also carry risks of misattribution and emotional dependence \citep{pnas2025anthro, maeda2024parasocial}. Recent studies have begun to explore how to measure these parasocial dynamics with AI \citep{omeish2025between, qi2025assistant} and understand their long-term psychological impacts \citep{payne2024ai, xie2022ai}.

The concept of an AI "persona" is increasingly relevant as developers craft specific personalities for their systems. The stability and perception of these personas are key research topics. Recent work has explored the use of LLMs in generating user personas for UX design \citep{hur2024utilizing}, as well as the potential for "generative diaries" to shape long-term human-robot interaction \citep{cho2023story}. Studies have evaluated how different modalities for presenting personas (e.g., chat vs. profile) affect user perception and behavior \citep{kaate2025evaluating, kaate2025fourth}. However, the challenge of maintaining a consistent and coherent persona across rapid model iterations remains a significant, yet underexplored, area of research. Our work addresses this gap by focusing on the "mutual wanting" between users and systems during such transitions, treating system constraints and user desires as equally important parts of the interaction design.

\section{Methodology}
Our study employed a two-pronged approach to investigate the "mutual wanting" between users and LLMs. We combined large-scale analysis of public discourse with controlled API probing to capture both user-generated reactions and system-level behaviors. 

\textbf{Data Collection and Analysis.} We collected a corpus of 1,748 Reddit posts from subreddits related to GPT models, focusing on a period surrounding a major version transition. The data was filtered to identify posts discussing perceived changes in the model's persona, such as warmth, creativity, and helpfulness. A subset of 300 posts was independently annotated by two researchers to identify distinct style clusters in user posts. This process achieved a Cohen's $\kappa$ of 0.795, indicating substantial agreement. We then extracted linguistic features from each post, including hedge rate, warmth markers, pronoun ratio, and token length. These features were used in a logistic regression model to predict post engagement (as measured by score).

\textbf{Behavioral Probing.} We developed a suite of controlled prompts to probe the behavior of different GPT model versions. These prompts were designed to elicit responses in situations with latent ambiguity, creative ideation, and emotional content. We measured metrics such as Confidence Disclosure Rate (CDR) and Structured Uncertainty Ratio (SUR) to quantify how models communicate uncertainty.

\textbf{Temporal Analysis.} To analyze style drift over time, we segmented our corpus into "pre" and "post" transition periods. We then compared the distribution of style clusters and the frequency of specific linguistic markers between these periods to identify significant shifts in user discourse and model behavior.

\section{Research Question and Narrative Arc}
	extbf{Core Question:} How do user desiderata and system-implied behavioral expectations co-evolve across major LLM version transitions, and where do tensions (“mutual wanting” misalignments) emerge that are observable in public discourse and controlled probe responses?

	extbf{Story Spine (Planned):}
\begin{enumerate}
  \item \emph{Context}: Rapid iteration of frontier LLMs produces perceived “persona drift.”
  \item \emph{Signal}: Users publicly articulate loss/gain framings (warmth, uncertainty style, creativity).
  \item \emph{Hypothesis}: These complaints cluster along stable tension axes (relational continuity vs optimization; openness vs guarded uncertainty; expressive richness vs efficiency; intimacy boundary negotiation).
  \item \emph{Measurement Plan}: (a) Authentic pre/post window corpus comparison; (b) Annotation of major relational / epistemic tags; (c) Regression + drift + probe triangulation; (d) Lexical–theme coupling to tie micro signals to macro axes.
  \item \emph{Outcome Vision}: A compact set of tension dashboards enabling proactive alignment observability for future model launches.
\end{enumerate}

	extbf{Significance}: Current evaluation paradigms under-weight relational stability. Our framework operationalizes lightweight, reproducible observability layers (annotation reliability, drift stability, probe deltas) without privileged internal telemetry. This supports transparency norms the community is calling for and can reduce rumor-driven trust erosion.

\section{Results (Status: Incomplete / No Claims Yet)}
	extbf{Transparency Notice.} The previously included tables (regression interactions, probe metrics, drift tokens) were automatically scaffolded but contain only headers or placeholder rows (no substantive analytical data). They have been removed to prevent inadvertent over-claiming. No inferential results are reported at this stage.

\subsection*{Current Empirical State}
\begin{itemize}
  \item \textbf{Corpus}: 1,748 recent Reddit posts (pilot-only slice). \emph{Missing}: Authenticated historical pre/post transition window dumps (blocking full drift analysis).
  \item \textbf{Annotation}: Pilot A/B CSVs exist (40+40 with 20 overlap). Agreement TSV shows preliminary Cohen's $\kappa=0.707$ (overlap=40) \citep{landis1977measurement}. Prior text incorrectly reported $0.795$—now retracted pending scale annotation.
  \item \textbf{Features}: Per-comment linguistic features + provisional style clusters (needs validation: small cluster instability; silhouette not yet documented in text).
  \item \textbf{Regression}: Model dataset limited to 60 pilot-labeled rows—insufficient for pre"+"post * transition interaction modeling; placeholder output not interpretable.
  \item \textbf{Probes}: Manifest generated; statistical aggregation script not yet executed for real effect estimates (no CIs, no hypothesis tests). Only metric header TSV present.
  \item \textbf{Drift}: No authentic pre/post aligned window pair normalized; drift token table empty; bootstrap + placebo tests unrun.
\end{itemize}

\subsection*{Planned Analytical Completion Roadmap}
\begin{enumerate}
  \item Acquire \& integrity-validate historical window dumps (two transitions if feasible) + normalize + coverage audit.
  \item Scale annotation to target N ($\geq 600$ combined) with overlap $\geq 15$\%; recompute overall \& per-tag $\kappa$; consensus merge.
  \item Rebuild regression with sufficient class support (MIN\_CLASS $\geq 8$) including interaction terms (pre/post * transition) + robustness (placebo shuffles, subreddit sensitivity).
  \item Execute lexical drift (log-odds with informative prior) + bootstrap stability + placebo.
  \item Run probe responses (two model generations) → compute CDR, SUR, ETD, SST, CPR, DRP + proportion tests.
  \item Couple lexical drift signals with regression interaction ORs + probe deltas for convergent tension evidence.
  \item Produce finalized figures: Fig2 (confusion heatmap post-scale), Fig3 (forest with real CIs), Fig4 (stable lexical drift), Fig5 (probe contrasts), Fig6 (lexical–theme coupling), Fig1 (system pipeline diagram now added below).
  \item Populate tables: Sampling/Coverage (updated), Enrichment Evaluation, Regression Interactions, Probe Metrics, Robustness \& Placebo, Drift Lexicon (appendix).
\end{enumerate}

\subsection*{Provisional Reliability Figure (Pilot)}
\begin{figure}[htbp]
  \centering
  \includegraphics[width=0.70\linewidth]{../pipeline/outputs/figs/fig2_confusion_heatmap}
  \caption{Pilot overlap confusion (counts normalized). \emph{Note}: Pilot-scale only; do not generalize.}
  \label{fig:confusion}
\end{figure}

\subsection*{System / Pipeline Overview}
\begin{figure}[htbp]
  \centering
  \includegraphics[width=0.80\linewidth]{../figures/drawio/Fig1}
  \caption{System + analysis pipeline schematic (minimal-resource observability stack).}
  \label{fig:system}
\end{figure}

\subsection*{Ethics and Data Integrity Guardrails (Planned Checks)}
Planned: (a) No synthetic data (enforced by raw source hash manifest); (b) Author hash salting documented; (c) Near-duplicate cap ($\leq 5$\%) with longest variant retention; (d) Parasocial risk metrics (ETD distribution) only reported in aggregate; (e) Placebo drift false-positive calibration.

\subsection*{Removed Prior Claims}
Any earlier numerical effect sizes (ORs, drift percentages, probe means) are withdrawn until full dataset assembly and validation. The manuscript will be updated once authentic analyses are complete.


\section{Discussion}
Our findings illuminate the delicate and dynamic balance in the "mutual wanting" between users and AI systems. The strong inter-annotator agreement on style clusters is a key result, suggesting that users exhibit consistent, recognizable patterns of expression when discussing AI models. These are not random complaints; they are structured responses. The regression analysis reinforces this, confirming that linguistic choices, such as hedging and verbosity, are not merely stylistic but are significantly associated with community engagement. This reflects a social evaluation process where the community collectively assesses the validity and nature of AI-related discourse. The very act of complaining about a model's "persona" is a social act, evaluated by peers.

The observed temporal drift in these style clusters, although modest, points to an evolving user response to the AI's changing behavior. The decrease in the dominant style cluster suggests a fragmentation or diversification of user reactions over time. This could reflect a growing sophistication in how users perceive and critique AI, moving beyond monolithic judgments to more nuanced assessments of specific traits. It may also indicate that as users adapt to a new model, their methods of expressing dissatisfaction change. This is a critical insight for developers: the user community is not a static entity but a dynamic one that co-evolves with the model.

The results from our behavioral probes provide a glimpse into the system's side of the "mutual wanting" equation. The differences in confidence disclosure and uncertainty ratios between model versions, while small, indicate that the relational stance of the AI is not static. These subtle shifts in how a model communicates its own uncertainty can have an outsized impact on user perception, contributing to the sense of "persona drift" that our study aims to capture. When a model that was previously cautious becomes more assertive, or vice versa, it violates user expectations and can erode trust.

This study has several limitations that future work should address. Our analysis is based on public Reddit data, which, while a rich source of spontaneous user feedback, may not be representative of the broader user base. The users who post on Reddit are likely more technically engaged and may have different expectations than more casual users. Furthermore, the correlational nature of our analysis does not allow for causal claims about the relationship between model changes and user reactions. Finally, our "pre" and "post" periods are defined around a single, albeit significant, version transition. The observed drift may not generalize to other types of model updates or different AI systems. Despite these limitations, our minimal-resource approach demonstrates that it is possible to gain significant insights into the human-AI relationship using publicly available data, providing a valuable starting point for more extensive research.

\section{Conclusion}
This study introduced a minimal-resource framework for analyzing the dynamic and often fraught relationship between users and evolving LLMs. By empirically grounding the concept of "mutual wanting," we have moved beyond simplistic notions of user satisfaction to a more nuanced understanding of the bidirectional expectations in human-AI interaction. Our analysis of public discourse and controlled behavioral probes reveals that the "persona" of an AI is not a static design feature but a co-constructed property that emerges from the interplay of user expression and system behavior. The measurable shifts we observed in both domains underscore the need for a relational approach to AI design, one that prioritizes transparency, consistency, and explicit management of user expectations during model updates.

Our findings have direct implications for the design and deployment of large-scale AI systems. The "move fast and break things" ethos of software development is ill-suited to the deployment of relational AI, where user trust and emotional investment are at stake. Instead, we advocate for a more deliberate and communicative approach to model transitions, one that includes clear articulation of changes, and mechanisms for user feedback.

Future work should expand this analysis across a wider range of model architectures, user populations, and cultural contexts. Longitudinal studies are needed to understand the long-term effects of "persona drift" on user trust and engagement. Furthermore, our "mutual wanting" framework could be extended to other domains of human-AI interaction, such as creative collaboration and personalized education. By continuing to map the complex terrain of human-AI relationships, we can work towards building AI systems that are not only more capable but also more trustworthy and aligned with human values.

\bibliographystyle{plainnat}
\bibliography{references}

% --- End main manuscript content ---

%%%%%%%%%%%%%%%%%%%%%%%%%%%%%%%%%%%%%%%%%%%%%%%%%%%%%%%%%%%%
\section*{Agents4Science AI Involvement Checklist}
\begin{enumerate}
  \item \textbf{Hypothesis development}: \\Answer: \involvementB{} \\Explanation: Core problem framing and mutual wanting construct seeded by human authors; AI assisted by rapidly enumerating tension axes, refining terminology, and surfacing adjacent literatures for validation.
  \item \textbf{Experimental design and implementation}: \\Answer: \involvementC{} \\Explanation: Pipeline code (fetching, sampling, annotation tooling, regression/drift scaffolds) was predominantly AI-drafted then human-audited for robustness, ethics, and feasibility constraints.
  \item \textbf{Analysis of data and interpretation of results}: \\Answer: \involvementB{} \\Explanation: Analytical plan (metrics, gating criteria, robustness checks) co-developed; humans set acceptance thresholds (e.g., kappa > 0.65, frequency floors) and AI generated alternative formulations and risk mitigations.
  \item \textbf{Writing}: \\Answer: \involvementC{} \\Explanation: Majority of prose (framework exposition, methodology detailing, robustness subsection) AI-generated with human curation for coherence, pruning template boilerplate, and ensuring accurate scoping (outline vs final claims).
  \item \textbf{Observed AI Limitations}: \\Description: Tended toward over-confident early interpretation of synthetic validation, required human intervention to clearly label provisional results; occasional citation hallucination risk mitigated via manual bib cross-check; needed guidance to avoid redundant robustness wording and to excise conference template filler.
\end{enumerate}

\newpage

\section*{Agents4Science Paper Checklist}
\begin{enumerate}
\item {\bf Claims} \\Answer: \answerYes{} \\Justification: Abstract and Introduction explicitly frame this as an outline/minimal-method plan without presenting final empirical findings (see Introduction, Planned Results Presentation).
\item {\bf Limitations} \\Answer: \answerYes{} \\Justification: Limitations are discussed in the ``Discussion'' section.
\item {\bf Theory assumptions and proofs} \\Answer: \answerNA{} \\Justification: No formal theorems or proofs are presented; work is empirical design and measurement specification.
\item {\bf Experimental result reproducibility} \\Answer: \answerYes{} \\Justification: Data pipeline components, gating criteria, sampling, and metrics are described in the ``Methodology'' section.
\item {\bf Open access to data and code} \\Answer: \answerNo{} \\Justification: Code will be released post-review. The abstract and introduction state the plan to provide reproducibility artifacts.
\item {\bf Experimental setting/details} \\Answer: \answerNA{} \\Justification: No final experiments yet; only planned metrics and pipeline steps described pre-data collection.
\item {\bf Experiment statistical significance} \\Answer: \answerNA{} \\Justification: No statistical result claims yet; future tests (e.g., drift log-odds z-scores) are planned but not executed.
\item {\bf Experiments compute resources} \\Answer: \answerNA{} \\Justification: Compute requirements minimal and future (Reddit fetch + lightweight regression); will detail upon execution.
\item {\bf Code of ethics} \\Answer: \answerYes{} \\Justification: Ethical considerations are mentioned in the abstract and introduction as ``ethical guardrails.''
\item {\bf Broader impacts} \\Answer: \answerYes{} \\Justification: Societal impacts are discussed in the ``Discussion'' and ``Conclusion'' sections.
\end{enumerate}

\end{document}
